\documentclass[10pt,a4paper]{amsart}
\usepackage[margin=19mm]{geometry}
\usepackage{amsmath,amssymb,mathtools}
\usepackage[scr=rsfs,cal=euler]{mathalfa}
\usepackage{xfrac}
\usepackage[unicode,hidelinks,bookmarks=false]{hyperref}
\newcommand{\ie}{i.e.\ }
\newcommand{\cf}{cf.\ }
\newcommand{\resp}{respectively\ }
\newcommand{\Subsection}{\subsection}
\providecommand{\nobreakdash}{\nobreak\hbox{-}}
\setlength{\emergencystretch}{2em}
\multlinegap0pt
\usepackage{amssymb}
\usepackage[all]{xy}
\newtheorem{lemma}{Lemma}
\newtheorem{theorem}{Theorem}
\newcommand{\LIM}{\mathbf{LIM}}
\def\c {{\mathfrak{c}}}
\title{Remarks on $d$-independent topological groups}
\author{Zhouxiang Huang, Dekui Peng,, Gao Zhang}
\date{Source: arXiv:2601.03000v1 (CC BY 4.0)}
\begin{document}
\maketitle

\noindent\emph{Source and licence.} Remarks on $d$-independent topological groups. Version arXiv:2601.03000v1 (CC BY 4.0), distributed by arXiv under the Creative Commons Attribution 4.0 International licence, http://creativecommons.org/licenses/by/4.0/, as recorded in the arXiv licence metadata for this exact version and shown on the abstract page https://arxiv.org/abs/2601.03000v1. Full licence notice: \texttt{licenses/arxiv-2601.03000-CC-BY-4.0.txt}.\par\smallskip

\noindent\textit{Editorial scope.} Counted unit: the article's theorem Th:eq (source0.tex lines 282--285). The article's complete original proof of it is delivered with it (source0.tex lines 287--340, one \texttt{proof} environment of 54 lines). No mathematics is written, repaired or re-derived: the statement and the proof are the article's own lines, in the article's own notation, and no retained line is substituted or re-flowed.  Retained uncounted as the source-local context the proof needs: lines 201--210; lines 240--247.  Nothing was omitted from any retained span, so the package carries no omission record.  No retained line contains a citation, so the wrapper carries no bibliography and no bibliography entry.  No retained line contains a figure environment, an image-inclusion command or drawing code, so no image is delivered and the package declares no asset dependency.  Editorial formatting only, in addition: an AMS \texttt{amsart} preamble that restates the article's own declarations and macros used by the retained lines, namely the environments \texttt{lemma}, \texttt{theorem} on separate counters, together with the standard packages those lines need.  The retained environments print in this document as Lemma 1, Theorem 1 in order of appearance, whereas the article prints the same items with the numbering of its own full document.  The complete native source of the article as distributed in the arXiv e-print for this exact version is bundled unchanged as \texttt{sources/192201/source0.tex}.
\begin{lemma}\label{indfam}
Let $G$ be an abelian group and let $H$ and $K$ be subgroups of $G$ such that
$H\cap K=\{0\}$.
Let $\mathcal{A}$ be an independent family of subgroups of $H$.
For each $A\in \mathcal{A}$, fix a homomorphism $f_A\colon A\to K$, and define
\[
B_A=\{a+f_A(a): a\in A\}.
\]
Then the family $\mathcal{B}=\{B_A: A\in \mathcal{A}\}$ is independent.
\end{lemma}

\begin{lemma}\label{Le:eq}
Let $G$ be a separable abelian topological group. Then the following are equivalent:
\begin{itemize}
\item[(1)] $G$ is $d$-independent;
\item[(2)] $G$ admits a family $\{D_\alpha:\alpha<\mathfrak c\}$ of independent countable dense subgroups;
\item[(3)] $G$ admits a family $\{D_\alpha:\alpha<\mathfrak c\}$ of pairwise independent countable dense subgroups.
\end{itemize}
\end{lemma}

\begin{theorem}\label{Th:eq}
A separable abelian group $G$ with a torsion-free $d$-independent subgroup
is itself $d$-independent.
\end{theorem}

\begin{proof}
Let $H$ be a torsion-free $d$-independent subgroup of $G$.
Since $G$ is separable, there exists a dense countable subgroup $K$ of $G$.
Set $S=K\cap H$, which is countable.

As $H$ is $d$-independent, we can choose a dense countable subgroup
$D_0$ of $H$ such that $S\cap D_0=\{0_G\}$.
By transfinite induction, for each ordinal $\alpha<\mathfrak{c}$ we may select
a dense countable subgroup $D_\alpha$ of $H$ such that
\[
D_\alpha\cap\Bigl(S+\sum_{\beta<\alpha}D_\beta\Bigr)=\{0_G\}.
\]
Hence the family
\[
\mathcal{A}:=\{D_\alpha:\alpha<\mathfrak{c}\}
\]
is independent and $A\cap K=\{0_G\}$, where $A=\bigoplus_{\alpha<\c} D_\alpha$.

Let $\mathbf{LIM}$ denote the set of all limit ordinals below $\mathfrak{c}$.
For each $\alpha<\mathfrak{c}$, choose a nonzero element
$b_\alpha\in D_\alpha$.
For every $\gamma\in\mathbf{LIM}$, define
\[
P_\gamma=\sum_{n\in\omega\setminus\{0\}}\langle b_{\gamma+n}\rangle.
\]
%Then, $P_\gamma$ is a free abelian group with $\omega$ many generators.
%
By the construction of the family $\{D_\alpha:\alpha<\mathfrak{c}\}$, this sum
is direct; hence $P_\gamma$ is a free abelian group of rank $\omega$.

Fix, for each $\gamma\in\mathbf{LIM}$, a surjective homomorphism
$f_\gamma\colon D_\gamma\oplus P_\gamma\to K$ such that
$D_\gamma\leq\ker f_\gamma$, and set
\[
Q_\gamma=\{x+f_\gamma(x):x\in D_\gamma\oplus P_\gamma\}.
\]

We claim that each $Q_\gamma$ is dense in $G$.
Indeed, $Q_\gamma\cap H$ is dense in $H$ since it contains $D_\gamma$.
Therefore $\overline{H}\leq\overline{Q_\gamma}$, and in particular
$D_\gamma\oplus P_\gamma\leq\overline{Q_\gamma}$.
Thus, for every $x\in D_\gamma\oplus P_\gamma$,
\[
f_\gamma(x)=(x+f_\gamma(x))-x\in Q_\gamma+\overline{Q_\gamma}
\subseteq\overline{Q_\gamma}.
\]
As $f_\gamma$ is surjective, it follows that $K\leq\overline{Q_\gamma}$.
Since $K$ is dense in $G$, we conclude that $Q_\gamma$ is dense in $G$.

The independence of the family
$\{Q_\gamma: \gamma\in \LIM\}$
follows directly from Lemma~\ref{indfam}, as $Q_\gamma\leq A$.
Therefore, by Lemma~\ref{Le:eq}, the group $G$ is $d$-independent.
\end{proof}

\end{document}
